
\section{Differential polynomials and the disappearance of subtractions}
\label{dpm:section}

The all-power expansion also supplies a complete calculus for a fixed
polynomial in digamma and its positive derivatives.  Three further
continuations \cite{dpm:reports} develop this direction.  We reconcile
here their common analytic core: derivatives remove the frozen residue,
so the derivative-containing part has an ordinary, unsubtracted global
pole expansion.  The integration-by-parts quotient, sixth/seventh
coordinate reductions and unequal-tail rank claims remain separate
questions; they are not consequences of this convergence theorem alone.

Let $P\in\mathbb C[u_0,u_1,\ldots,u_R]$ be a fixed polynomial, and put
\[
 f_P(z)=P\bigl(\psi(z),\psi'(z),\ldots,\psi^{(R)}(z)\bigr),\qquad
 \phi_P(T)=P(T,0,\ldots,0).
\]
Write $\mathcal J=(u_1,\ldots,u_R)$.  Thus $P\in\mathcal J$ means
that every monomial contains at least one positive derivative.
The coefficients of $P$ are fixed constants, with no dependence on $z$.
Choose a common pole-order bound $M$, and extract
\[
 p_j(n)=[\varepsilon^{-j}]f_P(-n+\varepsilon),\quad
 r(n)=p_1(n),\quad
 P_n(z)=\sum_{j=1}^M\frac{p_j(n)}{(z+n)^j}.
\]
Zero coefficients are understood.  The pole data come from the
convergent series \eqref{dgp:eq:pole-series}, differentiated a finite number
of times and multiplied.  They are finite harmonic polynomials.
For example a monomial $\prod_{a=0}^R u_a^{m_a}$ has pole order at
most $\sum_a(a+1)m_a$.

\begin{theorem}[Unsubtracted expansion for the derivative ideal]
\label{dpm:unsubtracted}
For $P\in\mathcal J$, the identity
\begin{equation}\label{dpm:ML}
 \boxed{f_P(z)=\sum_{n\ge0}P_n(z)}
\end{equation}
holds with absolute normal convergence away from the nonpositive
integers.  For some fixed $A\ge0$, as $n\to\infty$,
\begin{align}
 p_j(n)&=O\bigl((1+\log n)^A\bigr),&
 \Delta p_j(n)&=O\bigl(n^{-1}(1+\log n)^A\bigr),\label{dpm:coeff}\\
 r(n)&=O\bigl(n^{-1}(1+\log n)^A\bigr),&
 \Delta r(n)&=O\bigl(n^{-2}(1+\log n)^A\bigr).
 \label{dpm:residue}
\end{align}
Here $\Delta a(n)=a(n)-a(n-1)$.
\end{theorem}

\begin{proof}
Put $K(v)=-\pi\cot(\pi v)$.  Its differential equation is
$K'=K^2+\pi^2$.  Each positive derivative $K^{(a)}$ is therefore a
polynomial in $K$ divisible by $K^2+\pi^2$: differentiation preserves
that divisibility.  For any constant $L$, the polynomial
$P(L+K,K',\ldots)$ is divisible by $K^2+\pi^2$.  Divide by this
factor and integrate the resulting polynomial in $K$.  There is a
polynomial $T_L$ with
\[
 P(L+K(v),K'(v),\ldots)=\frac{d}{dv}T_L(K(v)).
\]
A meromorphic derivative has zero residue.  This is the exact frozen
residue cancellation; it uses the derivative in each monomial.

Fix $0<\rho<1/2$.  Reflection at $-n$ gives
\[
 \psi(-n+v)=K(v)+\log n+\eta_n(v),\qquad
 \eta_t(v)=\psi(t+1-v)-\log t\quad(t\ge1).
\]
Uniformly for $|v|\le\rho$, every fixed finite set of derivatives satisfies
\[
 \partial_v^a\eta_t(v)=O(t^{-1}),\qquad
 \partial_t\partial_v^a\eta_t(v)=O(t^{-2}).
\]
For $a=0$ this includes the subtraction of $\log t$ and its derivative;
for $a>0$ it follows directly from the positive-axis polygamma
expansions, which may also be obtained by differentiating the normally
convergent pole series.  Cauchy coefficient extraction on $|v|=\rho$
now bounds $p_j(t)$ and its $t$ derivative by the orders in
\eqref{dpm:coeff}.  For the residue, first subtract the frozen
polynomial with $L=\log t$.  Its residue is identically zero, while
the difference and its $t$ derivative have bounds
$O(t^{-1}(1+\log t)^A)$ and $O(t^{-2}(1+\log t)^A)$.
Integrating the $t$ derivative over $[n-1,n]$ proves
\eqref{dpm:residue}; finitely many small indices do not affect these
estimates.  No derivative of an unspecified remainder is used.

These bounds give absolute normal convergence of the proposed pole
series: the simple pole has the extra factor $n^{-1}$, and every
higher pole already has exponent at least two.  To determine the
entire remainder, use $|w|=R_N=N+1/2$.  Pole avoidance, reflection
and the same sectorial expansions give uniformly on this circle
\[
 |\psi(w)|\le C(1+\log R_N),\qquad
 |\psi^{(a)}(w)|\le C_a
       \bigl(R_N^{-1}+e^{-c|\operatorname{Im}w|}\bigr)
       \quad(a\ge1).
\]
The cotangent derivatives decay exponentially for large imaginary
part and remain bounded for small imaginary part because the circle
stays at least $1/2$ from each integer.  The right half-plane can
be estimated directly.  Thus, since every monomial has a positive
derivative,
\[
 |f_P(w)|\le C(1+\log R_N)^A
        \bigl(R_N^{-1}+e^{-c|\operatorname{Im}w|}\bigr).
\]
The angular integral of $e^{-cR_N|\sin\theta|}$ is $O(R_N^{-1})$.
Consequently, for fixed $z$,
\[
 \oint_{|w|=R_N}\frac{f_P(w)}{w-z}\,dw
       =O\bigl((1+\log R_N)^A/R_N\bigr)\longrightarrow0.
\]
The residue at $z$ is $f_P(z)$, and the one at $-n$ is $-P_n(z)$.
The residue theorem proves \eqref{dpm:ML}, including its normalization.
\end{proof}

\begin{theorem}[Ordinary finite-part coordinates for the derivative ideal]
\label{dpm:FP}
For $P\in\mathcal J$, set
\[
 J_P=\sum_{n\ge1}\sum_{j=2}^M
       \frac{\Delta p_j(n)}{(j-1)n^{j-1}},\qquad
 D_P(s)=\sum_{n\ge1}\frac{\Delta r(n)}{n^s}.
\]
Then $D_P$ is holomorphic for $\operatorname{Re}s>-1$, and the
original-coordinate finite part is
\begin{equation}\label{dpm:ordinary}
 \boxed{\operatorname{FP}_x\int_0^1 f_P(x)\,dx
   =J_P+D_P'(0)
   =J_P-\sum_{n\ge1}\Delta r(n)\log n.}
\end{equation}
Every displayed sum is ordinary and absolutely convergent.
\end{theorem}

\begin{proof}
After removing the finitely many pole terms at zero,
\eqref{dpm:ML} converges uniformly on $[0,1]$, so it may be integrated
there.  The simple pole at zero has finite part zero.  A higher pole
at zero contributes $-p_j(0)/(j-1)$.  A pole at $-n$, $n\ge1$,
contributes $r(n)\log(1+1/n)$ and
$p_j(n)(n^{1-j}-(n+1)^{1-j})/(j-1)$.
Use the finite summation identities
\eqref{dgp:eq:sbp-high} and \eqref{dgp:eq:sbp-simple}.
The initial higher-pole terms cancel the finite parts at zero.
All boundary terms vanish: $p_j(N)(N+1)^{1-j}\to0$ for $j\ge2$,
and $r(N)\log(N+1)\to0$ by \eqref{dpm:residue}.
This gives \eqref{dpm:ordinary}.  The same bound on $\Delta r$
proves normal convergence of $D_P$ and all its derivatives on compact
subsets of $\operatorname{Re}s>-1$.  It justifies the displayed
ordinary derivative at zero.
\end{proof}

\begin{corollary}[An arbitrary fixed differential polynomial]
\label{dpm:allP}
Write $P=\phi_P(u_0)+P_{\rm mix}$ with $P_{\rm mix}\in\mathcal J$.
The finite part of $f_P$ is the linear combination of the all-power
formulas for $\phi_P$ and \eqref{dpm:ordinary} for $P_{\rm mix}$.
In particular, the large-circle boundary polynomial and the frozen
residue polynomial depend only on $\phi_P$.  If
$\phi_P(T)=\sum a_kT^k$, they are respectively
\[
 C_P(L)=\sum_k\frac{a_k}{2\pi}\int_{-\pi}^{\pi}(L+i\theta)^k\,d\theta,
 \qquad
 R_P(L)=-\frac{\phi_P(L+i\pi)-\phi_P(L-i\pi)}{2i\pi}.
\]
The normalized global formula is
\[
 f_P(z)=P_0(z)+\kappa_P+
       \sum_{n\ge1}\bigl(P_n(z)-P_n(0)\bigr),
 \qquad \kappa_P=[\varepsilon^0]f_P(\varepsilon).
\]
\end{corollary}

\begin{proof}
Apply Theorems~\ref{dgp:thm:ML}, \ref{dgp:thm:contour} and
\ref{dgp:thm:all-power} to each scalar monomial, and the two preceding
theorems to $P_{\rm mix}$.  The scalar constant term contributes its
ordinary integral.  For the mixed part, \eqref{dpm:ML} implies
$\kappa_{P_{\rm mix}}=\sum_{n\ge1}(P_{\rm mix})_n(0)$ by taking the regular constant at zero;
this sum converges by \eqref{dpm:coeff}--\eqref{dpm:residue}.
Hence its unsubtracted expansion is also precisely the displayed
normalized expansion.  Linearity completes every assertion.
\end{proof}

\begin{corollary}[Anchored primitives for derivative-containing polynomials]
\label{dpm:primitive}
On $\mathbb C\setminus(-\infty,0]$, use the principal logarithm.
For $z,z_0$ in this domain,
\begin{align*}
 F_P(z;z_0)={}&\sum_{n\ge0}r(n)
       \bigl(\operatorname{Log}(z+n)-\operatorname{Log}(z_0+n)\bigr)\\
 &+\sum_{n\ge0}\sum_{j=2}^M\frac{p_j(n)}{j-1}
       \bigl((z_0+n)^{1-j}-(z+n)^{1-j}\bigr)
\end{align*}
converges normally, has $F_P(z_0;z_0)=0$ and
$\partial_zF_P(z;z_0)=f_P(z)$.  For complex anchors retain the difference
of logarithms; a principal logarithm of the ratio need not agree.
\end{corollary}

\begin{proof}
On a compact set, the logarithm difference is $O(n^{-1})$ and the
inverse-power difference in the $j$th term is $O(n^{-j})$.
The coefficient bounds imply normal convergence.  Differentiation
and \eqref{dpm:ML} prove the derivative and the anchor is immediate.
\end{proof}

As a normalization check, the derivative-containing polynomial
$P=u_0u_1$ gives
\[
 \operatorname{FP}_x\int_0^1\psi(x)\psi'(x)\,dx=\zeta(2).
\]
Indeed its primitive is $\psi^2/2$.  Its value at one is $\gamma^2/2$,
whereas its local constant at zero is $(\gamma^2-2\zeta(2))/2$.
More general total derivatives retain exactly this endpoint-minus-local-
constant rule.  Formula \eqref{dpm:ordinary} applies as well to mixed
polynomials which are not total derivatives.  It removes singular
spectral Laurent coordinates without asserting that their ordinary
convergent replacements reduce arithmetically to a smaller constant field.

For the nonlinear example $P=u_1^2$, write the local series as
$\psi'(-n+v)=v^{-2}+c_{2,n}+2c_{3,n}v+\cdots$.
The square has $p_4(n)=1$, $p_3(n)=0$, $p_2(n)=2c_{2,n}$ and
$r(n)=4c_{3,n}$.  Hence $\Delta p_2(n)=2n^{-2}$ and
$\Delta r(n)=4n^{-3}$.  Formula \eqref{dpm:ordinary} gives the
explicit ordinary identity
\[
 \operatorname{FP}_x\int_0^1\psi'(x)^2\,dx
       =2\zeta(3)+4\zeta'(3).
\]
The logarithmic sum is absolutely convergent.  This illustrates why a
positive derivative removes the singular spectral germ while retaining
an ordinary zeta derivative.  Independent local-subtraction quadrature
checks this value and the two total-derivative examples; it is a numerical
diagnostic, with no interval error claim.
